\subsection{Anillos completos de fracciones}
\label{subsection:0.7.6-fr}

\begin{env}[7.6.1]
\label{0.7.6.1-fr}
Sean $A$ un anillo linealmente topologizado, $(\mathfrak J_\lambda)$ un
sistema fundamental de vecindades de 0 en $A$ formado por ideales, $S$ una parte
multiplicativa de $A$. Sea $u_\lambda$ el homomorfismo canónico
$A\to A_\lambda=A/\mathfrak J_\lambda$, y para
$\mathfrak J_\mu\subset\mathfrak J_\lambda$, sea $u_{\lambda\mu}$
el homomorfismo canónico $A_\mu\to A_\lambda$. Pongamos
$S_\lambda=u_\lambda(S)$, de modo que
$u_{\lambda\mu}(S_\mu)=S_\lambda$. Los $u_{\lambda\mu}$ dan
canónicamente homomorfismos sobreyectivos
$S_\mu^{-1}A_\mu\to S_\lambda^{-1}A_\lambda$, para los cuales estos anillos
forman un sistema proyectivo; designemos por $A\{S^{-1}\}$ el límite
proyectivo de este sistema. Esta definición no depende del sistema fundamental
de vecindades $(\mathfrak J_\lambda)$ elegido; en efecto:
\end{env}

\begin{proposition}[7.6.2]
\label{0.7.6.2-fr}
El anillo $A\{S^{-1}\}$ es topológicamente isomorfo al separado completado del
anillo $S^{-1}A$ para la topología en la que un sistema fundamental de vecindades
de 0 está formado por los $S^{-1}\mathfrak J_\lambda$.
\end{proposition}

En efecto, si $v_\lambda$ es el homomorfismo canónico
$S^{-1}A\to S_\lambda^{-1}A_\lambda$ deducido de $u_\lambda$, el núcleo de
$v_\lambda$ es $S^{-1}\mathfrak J_\lambda$ y $v_\lambda$ es sobreyectivo,
de donde la proposición (7.2.1).

\begin{corollary}[7.6.3]
\label{0.7.6.3-fr}
Si $S'$ es la imagen canónica de $S$ en el separado completado $\widehat A$ de
$A$, $A\{S^{-1}\}$ se identifica canónicamente con
$\widehat A\{{S'}^{-1}\}$.
\end{corollary}

Se observará que aun si $A$ es separado y completo, no ocurre lo mismo con
$S^{-1}A$ para la topología definida por los
$S^{-1}\mathfrak J_\lambda$, como se ve, por ejemplo, tomando como $S$
el conjunto de los $f^n$ ($n\geq0$), donde $f$ es topológicamente nilpotente y no
nilpotente: en efecto, $S^{-1}A$ no se reduce a 0 y, por otra parte, para todo
$\lambda$ existe $n$ tal que $f^n\in J_\lambda$, por tanto
$1=f^n/f^n\in S^{-1}\mathfrak J_\lambda$ y
$S^{-1}\mathfrak J_\lambda=S^{-1}A$.

\oldpage[0\textsubscript{I}]{73}

\begin{corollary}[7.6.4]
\label{0.7.6.4-fr}
Si, en $A$, $0$ no es adherente a $S$, el anillo $A\{S^{-1}\}$ no se
reduce a $0$.
\end{corollary}

En efecto, $0$ no es adherente a $\{1\}$ en el anillo $S^{-1}A$; en caso contrario,
se tendría $1\in S^{-1}\mathfrak J_\lambda$ para todo ideal abierto
$\mathfrak J_\lambda$ de $A$, y resultaría que
$\mathfrak J_\lambda\cap S\neq\varnothing$ para todo $\lambda$, contrariamente
a la hipótesis.

\begin{env}[7.6.5]
\label{0.7.6.5-fr}
Diremos que $A\{S^{-1}\}$ es el \emph{anillo completo de las fracciones} de $A$
que tienen sus denominadores en $S$. Con las notaciones anteriores, es claro
que la imagen inversa de $S^{-1}\mathfrak J_\lambda$ en $A$ contiene
$\mathfrak J_\lambda$, por tanto la aplicación canónica $A\to S^{-1}A$ es
continua, y si se la compone con la aplicación canónica
$S^{-1}A\to A\{S^{-1}\}$, se obtiene un homomorfismo canónico continuo
$A\to A\{S^{-1}\}$, límite proyectivo de los homomorfismos
$A\to S_\lambda^{-1}A_\lambda$.
\end{env}

\begin{env}[7.6.6]
\label{0.7.6.6-fr}
El par formado por $A\{S^{-1}\}$ y por la aplicación canónica
$A\to A\{S^{-1}\}$ está caracterizado por la \emph{propiedad universal}
siguiente: todo homomorfismo continuo $u$ de $A$ en un anillo linealmente
topologizado $B$, separado y completo, tal que $u(S)$ está formado por elementos
invertibles en $B$, se factoriza de una sola manera en
$A\to A\{S^{-1}\}\xrightarrow{u'}B$, donde $u'$ es continuo. En efecto, $u$ se
factoriza de una sola manera en $A\to S^{-1}A\xrightarrow{v'}B$; como para
todo ideal abierto $\mathfrak K$ de $B$, $u^{-1}(\mathfrak K)$ contiene un
$\mathfrak J_\lambda$, ${v'}^{-1}(\mathfrak K)$ contiene necesariamente
$S^{-1}\mathfrak J_\lambda$, por tanto $v'$ es continuo; puesto que $B$ es separado y
completo, $v'$ se factoriza de una sola manera en
$S^{-1}A\to A\{S^{-1}\}\xrightarrow{u'}B$, donde $u'$ es continuo; de donde nuestra
afirmación.
\end{env}

\begin{env}[7.6.7]
\label{0.7.6.7-fr}
Sean $B$ un segundo anillo linealmente topologizado, $T$ una parte
multiplicativa de $B$, $\varphi:A\to B$ un homomorfismo continuo tal que
$\varphi(S)\subset T$. Según lo que precede, el homomorfismo continuo
$A\xrightarrow{\varphi}B\to B\{T^{-1}\}$ se factoriza de manera única en
$A\to A\{S^{-1}\}\xrightarrow{\varphi'}B\{T^{-1}\}$, donde $\varphi'$ es
continuo. En particular, si $B=A$ y si $\varphi$ es la identidad, se ve que
para $S\subset T$ se tiene un homomorfismo continuo
$\rho^{T,S}:A\{S^{-1}\}\to A\{T^{-1}\}$ obtenido por paso al separado completado
a partir de $S^{-1}A\to T^{-1}A$; si $U$ es una tercera parte
multiplicativa de $A$ tal que $S\subset T\subset U$, se tiene
$\rho^{U,S}=\rho^{U,T}\circ\rho^{T,S}$.
\end{env}

\begin{env}[7.6.8]
\label{0.7.6.8-fr}
Sean $S_1$, $S_2$ dos partes multiplicativas de $A$, y sea $S_2'$ la imagen
canónica de $S_2$ en $A\{S_1^{-1}\}$; se tiene entonces un isomorfismo
topológico canónico
$A\{(S_1S_2)^{-1}\}\xrightarrow{\sim}A\{S_1^{-1}\}\{{S_2'}^{-1}\}$,
como se ve partiendo del isomorfismo canónico
$(S_1S_2)^{-1}A\xrightarrow{\sim}{S_2''}^{-1}(S_1^{-1}A)$ (donde $S_2''$ es
la imagen canónica de $S_2$ en $S_1^{-1}A$), que es bicontinuo.
\end{env}

\begin{env}[7.6.9]
\label{0.7.6.9-fr}
Sea $\mathfrak a$ un ideal \emph{abierto} de $A$; se puede suponer que
$\mathfrak J_\lambda\subset\mathfrak a$ para todo $\lambda$, y, por consiguiente,
$S^{-1}\mathfrak J_\lambda\subset S^{-1}\mathfrak a$ en el anillo
$S^{-1}A$, es decir, $S^{-1}\mathfrak a$ es un ideal \emph{abierto} de
$S^{-1}A$; designaremos por $\mathfrak a\{S^{-1}\}$ su separado completado,
igual a $\varprojlim(S^{-1}\mathfrak a/S^{-1}\mathfrak J_\lambda)$, que es un
ideal \emph{abierto} de $A\{S^{-1}\}$, isomorfo a la adherencia de la imagen
canónica de $S^{-1}\mathfrak a$. Además, \emph{el anillo discreto
$A\{S^{-1}\}/\mathfrak a\{S^{-1}\}$ es canónicamente isomorfo a
$S^{-1}A/S^{-1}\mathfrak a=S^{-1}(A/\mathfrak a)$}. Recíprocamente, si
$\mathfrak a'$ es un ideal abierto de $A\{S^{-1}\}$, $\mathfrak a'$ contiene
un ideal de la forma $\mathfrak J_\lambda\{S^{-1}\}$, por tanto es la imagen
inversa de un ideal de $S^{-1}A/S^{-1}\mathfrak J_\lambda$, que es
necesariamente (1.2.6) de la forma $S^{-1}\mathfrak a$, donde
$\mathfrak a\supset J_\lambda$. Se concluye que se tiene
$\mathfrak a'=\mathfrak a\{S^{-1}\}$. En particular (1.2.6):
\end{env}

\begin{proposition}[7.6.10]
\label{0.7.6.10-fr}
La aplicación $\mathfrak p\to\mathfrak p\{S^{-1}\}$ es una biyección
creciente del conjunto de los ideales \textup{primos abiertos} $\mathfrak p$ de $A$ tales
que $\mathfrak p\cap S=\varnothing$ sobre el conjunto de los ideales \textup{primos abiertos}
\oldpage[0\textsubscript{I}]{74}
de $A\{S^{-1}\}$; además, el cuerpo de fracciones de
$A\{S^{-1}\}/\mathfrak p\{S^{-1}\}$ es canónicamente isomorfo al de
$A/\mathfrak p$.
\end{proposition}

\begin{proposition}[7.6.11]
\label{0.7.6.11-fr}
\begin{enumerate}
  \item[\rm(i)] Si $A$ es un anillo admisible, lo mismo ocurre con
    $A'=A\{S^{-1}\}$ y para todo ideal de definición $\mathfrak J$ de $A$,
    $\mathfrak J'=\mathfrak J\{S^{-1}\}$ es un ideal de definición de $A'$.
  \item[\rm(ii)] Sean $A$ un anillo ádico, $\mathfrak J$ un ideal de
    definición de $A$ tal que $\mathfrak J/\mathfrak J^2$ es de tipo finito
    sobre $A/\mathfrak J$; entonces $A'$ es un anillo $\mathfrak J'$-ádico y
    $\mathfrak J'/\mathfrak J'^2$ es de tipo finito sobre
    $A'/\mathfrak J'$. Si además $A$ es noetheriano, lo mismo ocurre con
    $A'$.
\end{enumerate}
\end{proposition}

\begin{enumerate}
  \item[\rm(i)] Si $\mathfrak J$ es un ideal de definición en $A$, es
    claro que $S^{-1}\mathfrak J$ es un ideal de definición en el anillo
    topológico $S^{-1}A$, pues se tiene
    $(S^{-1}\mathfrak J)^n=S^{-1}\mathfrak J^n$. Sea $A''$ el anillo separado
    asociado a $S^{-1}A$, $\mathfrak J''$ la imagen de $S^{-1}\mathfrak J$ en
    $A''$; la imagen de $S^{-1}\mathfrak J^n$ es $\mathfrak J''^n$, por tanto
    $\mathfrak J''^n$ tiende a 0 en $A''$; como $\mathfrak J'$ es
    la adherencia de $\mathfrak J''$ en $A'$, $\mathfrak J'^n$ está contenido
    en la adherencia de $\mathfrak J''^n$, por tanto tiende a 0 en $A'$.
  \item[\rm(ii)] Pongamos $A_i=A/\mathfrak J^{i+1}$, y para $i\leq j$, sea
    $u_{ij}$ el homomorfismo canónico
    $A/\mathfrak J^{j+1}\to A/\mathfrak J^{i+1}$; sea $S_i$ la imagen
    canónica de $S$ en $A_i$, y pongamos $A_i'=S_i^{-1}A_i$; sea finalmente
    $u_{ij}':A_j'\to A_i'$ el homomorfismo deducido canónicamente de $u_{ij}$.
    Mostremos que el sistema proyectivo $(A_i',u_{ij}')$ verifica las condiciones
    de la prop.~(7.2.7): es claro que los $u_{ij}'$ son sobreyectivos;
    por otra parte, el núcleo de $u_{ij}'$ es
    $S_j^{-1}(\mathfrak J^{i+1}/\mathfrak J^{j+1})$ (1.3.2), igual a
    $\mathfrak J_j'^{i+1}$, donde
    $\mathfrak J_j'=S_j^{-1}(\mathfrak J/\mathfrak J^{j+1})$; finalmente,
    $\mathfrak J_1'/\mathfrak J_1'^2=S_1^{-1}(\mathfrak J/\mathfrak J^2)$,
    y como $\mathfrak J/\mathfrak J^2$ es de tipo finito sobre
    $A/\mathfrak J^2$, $\mathfrak J_1'/\mathfrak J_1'^2$ es de tipo finito
    sobre $A_1'$. Finalmente, si $A$ es noetheriano, lo mismo ocurre con
    $A_0'=S_0^{-1}(A/\mathfrak J)$, lo que termina de demostrar la proposición
    (7.2.8).
\end{enumerate}

\begin{corollary}[7.6.12]
\label{0.7.6.12-fr}
Bajo las hipótesis de \textup{(7.6.11, (ii))}, se tiene
$(\mathfrak J\{S^{-1}\})^n=\mathfrak J^n\{S^{-1}\}$.
\end{corollary}

Esto resulta, en efecto, de (7.2.7) y de la demostración de (7.6.11).

\begin{proposition}[7.6.13]
\label{0.7.6.13-fr}
Sean $A$ un anillo ádico noetheriano, $S$ una parte multiplicativa de $A$;
entonces $A\{S^{-1}\}$ es un $A$-módulo plano.
\end{proposition}

En efecto, si $\mathfrak J$ es un ideal de definición de $A$,
$A\{S^{-1}\}$ es el separado completado del anillo noetheriano $S^{-1}A$ provisto
de la topología $S^{-1}\mathfrak J$-preádica; por consiguiente (7.3.3),
$A\{S^{-1}\}$ es un $S^{-1}A$-módulo plano; como $S^{-1}A$ es un
$A$-módulo plano (6.3.1), la proposición resulta de la transitividad de la
planitud (6.2.1).

\begin{corollary}[7.6.14]
\label{0.7.6.14-fr}
Bajo las hipótesis de \textup{(7.6.13)}, sea $S'\subset S$ una segunda parte
multiplicativa de $A$; entonces $A\{S^{-1}\}$ es un
$A\{{S'}^{-1}\}$-módulo plano.
\end{corollary}

En efecto (7.6.8), $A\{S^{-1}\}$ se identifica canónicamente con
$A\{{S'}^{-1}\}\{S_0^{-1}\}$, donde $S_0$ es la imagen canónica de $S$ en
$A\{{S'}^{-1}\}$, y $A\{{S'}^{-1}\}$ es noetheriano (7.6.11).

\begin{env}[7.6.15]
\label{0.7.6.15-fr}
Para todo elemento $f$ de un anillo linealmente topologizado $A$,
designaremos por $A_{\{f\}}$ el anillo completo de fracciones $A\{S_f^{-1}\}$,
donde $S_f$ es el conjunto multiplicativo de los $f^n$ ($n\geq 0$); para todo
ideal abierto $\mathfrak a$ de $A$, escribiremos $\mathfrak a_{\{f\}}$ en
lugar de $\mathfrak a\{S_f^{-1}\}$. Si $g$ es un segundo elemento de $A$, se tiene
un homomorfismo canónico continuo $A_{\{f\}}\to A_{\{fg\}}$ (7.6.7).
Cuando $f$ recorre una parte multiplicativa $S$ de $A$, los $A_{\{f\}}$
forman, pues, un sistema inductivo filtrante de anillos para los homomorfismos
anteriores; pondremos
$A_{\{S\}}=\varinjlim_{f\in S}A_{\{f\}}$. Para todo $f\in S$, se tiene un
homomorfismo $A_{\{f\}}\to A\{S^{-1}\}$ (7.6.7), y
\oldpage[0\textsubscript{I}]{75}
estos homomorfismos forman un sistema inductivo; por paso al límite
inductivo, definen, por tanto, un homomorfismo canónico
$A_{\{S\}}\to A\{S^{-1}\}$.
\end{env}

\begin{proposition}[7.6.16]
\label{0.7.6.16-fr}
Si $A$ es un anillo noetheriano, $A\{S^{-1}\}$ es un módulo plano sobre
$A_{\{S\}}$.
\end{proposition}

En efecto (7.6.14), $A\{S^{-1}\}$ es plano sobre cada uno de los anillos
$A_{\{f\}}$ para $f\in S$, y la conclusión resulta de (6.2.3).

\begin{proposition}[7.6.17]
\label{0.7.6.17-fr}
Sea $\mathfrak p$ un ideal primo abierto en un anillo admisible $A$, y
sea $S=A\setmin\mathfrak p$. Entonces los anillos $A\{S^{-1}\}$ y $A_{\{S\}}$ son
anillos locales, el homomorfismo canónico
$A_{\{S\}}\to A\{S^{-1}\}$ es local y los cuerpos residuales de $A_{\{S\}}$
y $A\{S^{-1}\}$ son canónicamente isomorfos al cuerpo de fracciones de
$A/\mathfrak p$.
\end{proposition}

En efecto, sea $\mathfrak J\subset\mathfrak p$ un ideal de definición de
$A$; se tiene $S^{-1}\mathfrak J\subset S^{-1}\mathfrak p=\mathfrak p
A_{\mathfrak p}$, por tanto $A_{\mathfrak p}/S^{-1}\mathfrak J$ es un anillo
local; se concluye de (7.1.12), (7.6.9) y (7.6.11, (i)) que
$A\{S^{-1}\}$ es un anillo local. Pongamos
$\mathfrak m=\varinjlim_{f\in S}\mathfrak p_{\{f\}}$, que es un ideal de
$A_{\{S\}}$; vamos a ver que todo elemento de $A_{\{S\}}$ que no pertenece
a $\mathfrak m$ es invertible. En efecto, tal elemento es la imagen en
$A_{\{S\}}$ de un elemento $z\in A_{\{f\}}$ que no pertenece a
$\mathfrak p_{\{f\}}$, para un $f\in S$; su imagen canónica $z_0$ en
$A_{\{f\}}/\mathfrak J_{\{f\}}=S_f^{-1}(A/\mathfrak J)$ no pertenece, pues,
a $S_f^{-1}(\mathfrak p/\mathfrak J)$ (7.6.9), lo que significa que
$z_0=\overline{x}/\overline{f}^k$, donde $x\notin\mathfrak p$ y
$\overline{x},\overline{f}$ son las clases de $x,f$ mod.~$\mathfrak J$.
Como $x\in S$, se tiene $g=xf\in S$, y en $S_g^{-1}A$,
$y_0=x^{k+1}/g^k$, imagen canónica de $x/f^k\in S_f^{-1}A$, admite un
inverso $x^{k-1}f^{2k}/g^k$. Esto implica \emph{a fortiori} que la imagen de
$y_0$ en $S_g^{-1}A/S_g^{-1}\mathfrak J$ es invertible, por tanto (7.6.9 y
7.1.12) la imagen canónica $y$ de $z$ en $A_{\{g\}}$ es invertible;
la imagen de $z$ en $A_{\{S\}}$ (igual a la de $y$) es, por consiguiente,
invertible. Se ve, pues, que $A_{\{S\}}$ es un anillo local de ideal maximal
$\mathfrak m$; además, la imagen de $\mathfrak p_{\{f\}}$ en
$A\{S^{-1}\}$ está contenida en el ideal maximal $\mathfrak p\{S^{-1}\}$ de
este anillo; \emph{a fortiori}, la imagen de $\mathfrak m$ en
$A\{S^{-1}\}$ está contenida en $\mathfrak p\{S^{-1}\}$, por tanto
el homomorfismo canónico $A_{\{S\}}\to A\{S^{-1}\}$ es local. Finalmente,
como todo elemento de $A\{S^{-1}\}/\mathfrak p\{S^{-1}\}$ es imagen de un
elemento de un anillo $S_f^{-1}A$ para un $f\in S$ adecuado,
el homomorfismo $A_{\{S\}}\to A\{S^{-1}\}/\mathfrak p\{S^{-1}\}$ es
sobreyectivo, y da, por tanto, por paso a los cocientes un isomorfismo de los cuerpos
residuales.

\begin{corollary}[7.6.18]
\label{0.7.6.18-fr}
Bajo las hipótesis de \textup{(7.6.17)}, si se supone además que $A$ es un anillo
ádico noetheriano, los anillos locales $A\{S^{-1}\}$ y $A_{\{S\}}$ son
noetherianos, y $A\{S^{-1}\}$ es un $A_{\{S\}}$-módulo fielmente plano.
\end{corollary}

Se sabe ya (7.6.11, (ii)) que $A\{S^{-1}\}$ es noetheriano y
$A_{\{S\}}$-plano (7.6.16); como el homomorfismo
$A_{\{S\}}\to A\{S^{-1}\}$ es local, se concluye que $A\{S^{-1}\}$ es
un $A_{\{S\}}$-módulo fielmente plano (6.6.2), y, por consiguiente, que $A_{\{S\}}$
es noetheriano (6.5.2).

